% ============================================================
%  SOAL + PEMBAHASAN  --  MATEMATIKA  --  SSU-ITB 2026
%  INTEN Set 1--6  (nomor 1--60) | Paket 3
% ============================================================

\documentclass[11pt]{article}

\usepackage[utf8]{inputenc}
\usepackage[T1]{fontenc}
\usepackage[a4paper,margin=2.5cm,top=2.8cm]{geometry}
\usepackage{amsmath,amssymb}
\usepackage{graphicx}
\usepackage{enumitem}
\usepackage{multicol}
\usepackage{xcolor}
\usepackage[most]{tcolorbox}
\usepackage{fancyhdr}
\usepackage{booktabs}
\usepackage{icomma}
\usepackage{array}

% --- Warna Matematika ---
\definecolor{mapelcolor}{RGB}{20,110,60}
\definecolor{mapelbg}{RGB}{228,246,234}
\definecolor{pembcolor}{RGB}{100,50,130}
\definecolor{pembbg}{RGB}{245,238,252}
\definecolor{gold}{RGB}{210,160,30}

\pagestyle{fancy}\fancyhf{}
\lhead{\small SSU-ITB 2026 --- Pembahasan}
\rhead{\small Paket 3}
\cfoot{\small\thepage}
\renewcommand{\headrulewidth}{0.4pt}

\newtcolorbox{soal}[1]{
  breakable,
  colback=mapelbg, colframe=mapelcolor,
  coltitle=white, fonttitle=\bfseries\sffamily,
  title=Nomor #1,
  boxrule=0.8pt, arc=3pt,
  left=3mm, right=3mm, top=2mm, bottom=2mm}

\newtcolorbox{pembox}{
  breakable,
  colback=pembbg, colframe=pembcolor,
  coltitle=white, fonttitle=\bfseries\sffamily,
  title=Pembahasan,
  boxrule=0.8pt, arc=3pt,
  left=3mm, right=3mm, top=2mm, bottom=2mm}

\newcommand{\jawaban}[1]{%
  \noindent\colorbox{gold!30}{\parbox{\dimexpr\linewidth-2\fboxsep\relax}{%
    \textbf{Jawaban:} #1}}\par\smallskip}

\newcommand{\konsep}{\noindent\textit{Konsep:}\ \ignorespaces}

\newenvironment{pil}{%
  \begin{enumerate}[label=(\Alph*),leftmargin=*,itemsep=1pt,topsep=3pt]}%
  {\end{enumerate}}
\newenvironment{pilB}{%
  \par\smallskip\begin{multicols}{2}
  \begin{enumerate}[label=(\Alph*),leftmargin=*,itemsep=1pt,topsep=3pt]}%
  {\end{enumerate}\end{multicols}}

\newcommand{\gambar}[2][0.52\textwidth]{%
  \begin{center}\includegraphics[width=#1]{#2}\end{center}}

\linespread{1.05}

% ===================================================================
\begin{document}
\thispagestyle{empty}

\begin{center}
  \fbox{\parbox{0.6\textwidth}{\centering\bfseries\large
    LEMBAR SOAL DAN PEMBAHASAN}}
\end{center}
\vspace{0.5cm}

\begin{center}
  {\LARGE\bfseries MATEMATIKA}
\end{center}
\vspace{0.5cm}

\begin{center}
  \begin{tabular}{@{\hspace{2cm}}ll@{\hspace{0.4cm}:@{\hspace{0.4cm}}}l}
    Jenjang      & & SMA / Sederajat               \\[0.3em]
    Hari/Tanggal & & \underline{\hspace{5cm}}      \\[0.3em]
    Waktu        & & \underline{\hspace{5cm}}      \\[0.3em]
    Paket Soal   & & 3                             \\
  \end{tabular}
\end{center}

\vspace{0.6cm}
\noindent\rule{\linewidth}{0.6pt}
\vspace{0.3cm}

\noindent\textbf{\underline{PETUNJUK UMUM}}
\begin{enumerate}[leftmargin=*,itemsep=4pt,topsep=4pt]
  \item Anda \textbf{tidak} diperbolehkan menggunakan sumber-sumber eksternal,
        termasuk internet, \textit{large language model} seperti ChatGPT, Claude,
        LLaMA, Gemini, Meta AI, dan sebagainya. Anda sangat dianjurkan untuk
        mencoba mengerjakan sendiri terlebih dahulu karena evaluasi ini dibuat
        untuk \textbf{mengukur} tingkat kepahaman; nilai $<70$ mengindikasikan
        \textbf{tidak lulus}, dan jika sudah melewati \textit{deadline} maka
        dianggap tidak mengerjakan yang berarti \textbf{tidak lulus}.

  \item Dokumen ini merupakan \textbf{lembar soal sekaligus pembahasan}; gunakan
        pada tahap \textbf{Review}. Di tahap ini Anda diharapkan mengerjakan ulang
        dengan disertakan penjelasan/pembelaan mengapa revisi jawaban adalah pilihan
        tersebut, dan tetap mencantumkan penjelasan pada soal yang walaupun pada
        penilaian sudah benar.

  \item Anda dianjurkan untuk mencantumkan caranya walaupun tidak termasuk
        penilaian, mengingat sifatnya adalah sebagai bahan pembelajaran.
        Mohon bertanggung jawab atas pekerjaan Anda sendiri.

  \item Terdapat tiga jenis soal: \textbf{Pilihan Ganda Biasa} (5 opsi, pilih 1);
        \textbf{Pilihan Ganda Majemuk Kompleks} (tabel \textbf{Benar}/\textbf{Salah},
        tiap pernyataan dinilai mandiri); dan \textbf{Isian Singkat}
        (tulis nilai/ekspresi).

  \item Soal berjumlah \textbf{60 butir}: nomor 1--20 dari Set 1--2 latihan INTEN
        SSU-ITB 2026, nomor 21--40 dari Set 3--4, dan nomor 41--60 dari Set 5--6.

  \item Isilah detail informasi berikut.\\[0.3em]
        \textbf{Nama}\hspace{0.45cm};\enspace\underline{\hspace{8cm}}
\end{enumerate}

\vspace{0.6cm}
\noindent\textit{Dengan menandatangani kolom di bawah ini, saya menyatakan telah membaca,
memahami, dan menyetujui seluruh peraturan yang tercantum di atas.}

\vspace{0.6cm}
\noindent\fbox{\parbox{5cm}{\vspace{2.2cm}\centering\small Tanda Tangan Peserta\vspace{0.3cm}}}

\clearpage

% ===================================================================
%  SET 1  (nomor 1--10)
% ===================================================================

\begin{soal}{1}
Jika $9^{1/6}\cdot x=3$, maka nilai $x$ adalah \ldots
\begin{pilB}
  \item $3^{4/6}$ \item $3^{4/36}$ \item $3^{6/2}$ \item $3^{6}$ \item $3^{1/3}$
\end{pilB}
\end{soal}
\begin{pembox}
\jawaban{(A) $x=3^{4/6}$}
\konsep $(a^m)^n=a^{mn}$ dan $a^m\cdot a^n=a^{m+n}$.

$9^{1/6}=(3^2)^{1/6}=3^{1/3}$. Dari $3^{1/3}\cdot x=3^1$:
$x=3^{1-1/3}=3^{2/3}=3^{4/6}$.
\end{pembox}

\begin{soal}{2}
Jika titik $A(4,15)$ terletak pada grafik fungsi $g(x)=\dfrac{2x+7}{x^2-b}$,
maka nilai $b$ adalah \ldots
\begin{pilB}
  \item $16$ \item $14$ \item $13$ \item $15$ \item $12$
\end{pilB}
\end{soal}
\begin{pembox}
\jawaban{(D) $b=15$}
$g(4)=15\;\Rightarrow\;15=\dfrac{2(4)+7}{16-b}=\dfrac{15}{16-b}\;\Rightarrow\;16-b=1\;\Rightarrow\;b=15$.
\end{pembox}

\begin{soal}{3}
Diketahui $x_0,\ldots,x_{100}$ adalah barisan aritmetika, $x_0=3$, $x_{100}=9$.
Tentukan kebenaran setiap pernyataan berikut!

\par\smallskip
{\renewcommand{\arraystretch}{1.4}%
\begin{tabular}{@{}p{0.76\linewidth}cc@{}}
\hline\textbf{Pernyataan}&\textbf{Benar}&\textbf{Salah}\\\hline
(1)~Beda $= \dfrac{3}{50}$.&$\bigcirc$&$\bigcirc$\\
(2)~$x_{50}=6$.&$\bigcirc$&$\bigcirc$\\
(3)~$x_{25}=4{,}5$.&$\bigcirc$&$\bigcirc$\\
(4)~$x_{75}=8$.&$\bigcirc$&$\bigcirc$\\\hline
\end{tabular}}
\end{soal}
\begin{pembox}
\jawaban{(1)~\textbf{B}\quad(2)~\textbf{B}\quad(3)~\textbf{B}\quad(4)~\textbf{S}}
$d=\dfrac{9-3}{100}=\dfrac{3}{50}$.\quad$x_n=3+n\cdot\dfrac{3}{50}$.
\begin{enumerate}[leftmargin=*,topsep=2pt,itemsep=2pt]
\item \textbf{Benar.} $d=\tfrac{3}{50}$.
\item \textbf{Benar.} $x_{50}=3+50\cdot\tfrac{3}{50}=6$ (juga $\tfrac{3+9}{2}=6$).
\item \textbf{Benar.} $x_{25}=3+\tfrac{3}{2}=4{,}5$.
\item \textbf{Salah.} $x_{75}=3+\tfrac{9}{2}=7{,}5\neq8$.
\end{enumerate}
\end{pembox}

\begin{soal}{4}
Diberikan $f(x)=\dfrac{5x^2-9x+3}{(x-1)^6}$.
Tentukan kebenaran setiap pernyataan berikut!

\par\smallskip
{\renewcommand{\arraystretch}{1.4}%
\begin{tabular}{@{}p{0.76\linewidth}cc@{}}
\hline\textbf{Pernyataan}&\textbf{Benar}&\textbf{Salah}\\\hline
(1)~$f'(x)$ terdefinisi untuk semua $x\neq1$.&$\bigcirc$&$\bigcirc$\\
(2)~$f'(2)=-19$.&$\bigcirc$&$\bigcirc$\\
(3)~$f'(2)>0$.&$\bigcirc$&$\bigcirc$\\
(4)~Jika $g(x)=(x-1)^6f(x)$, maka $g(2)=5$.&$\bigcirc$&$\bigcirc$\\\hline
\end{tabular}}
\end{soal}
\begin{pembox}
\jawaban{(1)~\textbf{B}\quad(2)~\textbf{B}\quad(3)~\textbf{S}\quad(4)~\textbf{B}}
$f'(x)=(10x-9)(x-1)^{-6}-6(5x^2-9x+3)(x-1)^{-7}$.
\begin{enumerate}[leftmargin=*,topsep=2pt,itemsep=2pt]
\item \textbf{Benar.} $f'$ hanya tidak ada saat $x=1$.
\item \textbf{Benar.} $f'(2)=11(1)-6(5)(1)=11-30=-19$.
\item \textbf{Salah.} $f'(2)=-19<0$.
\item \textbf{Benar.} $g(2)=(1)^6\cdot f(2)=20-18+3=5$.
\end{enumerate}
\end{pembox}

\begin{soal}{5}
Jika $x^3-2x+5$ dibagi oleh $1-x$, maka sisanya adalah \ldots
\begin{pilB}\item $4$\item $5$\item $6$\item $7$\item $8$\end{pilB}
\end{soal}
\begin{pembox}
\jawaban{(A) Sisa $= 4$}
\konsep Teorema Sisa: sisa pembagian $P(x)$ oleh $x-a$ adalah $P(a)$.

Pembagi $1-x$ berakar $x=1$: $P(1)=1-2+5=4$.
\end{pembox}

\begin{soal}{6}
Diketahui $f(x)=x^3+(a-1)x^2-2$. Jika $f(-2)=26$, maka nilai $a$ adalah \ldots
\begin{pilB}\item $5$\item $10$\item $15$\item $20$\item $25$\end{pilB}
\end{soal}
\begin{pembox}
\jawaban{(B) $a=10$}
$f(-2)=-8+4(a-1)-2=4a-14=26\;\Rightarrow\;4a=40\;\Rightarrow\;a=10$.
\end{pembox}

\begin{soal}{7}
Tiga buah dadu enam sisi dilempar bersamaan.
Tentukan kebenaran setiap pernyataan berikut!

\par\smallskip
{\renewcommand{\arraystretch}{1.4}%
\begin{tabular}{@{}p{0.76\linewidth}cc@{}}
\hline\textbf{Pernyataan}&\textbf{Benar}&\textbf{Salah}\\\hline
(1)~Ruang sampel berjumlah $216$.&$\bigcirc$&$\bigcirc$\\
(2)~Banyaknya cara agar jumlah ketiga mata dadu $=7$ adalah $15$.&$\bigcirc$&$\bigcirc$\\
(3)~Peluang jumlah $=7$ adalah $\dfrac{15}{216}$.&$\bigcirc$&$\bigcirc$\\
(4)~Peluang jumlah $=7$ adalah $\dfrac{1}{12}$.&$\bigcirc$&$\bigcirc$\\\hline
\end{tabular}}
\end{soal}
\begin{pembox}
\jawaban{(1)~\textbf{B}\quad(2)~\textbf{B}\quad(3)~\textbf{B}\quad(4)~\textbf{S}}
\begin{enumerate}[leftmargin=*,topsep=2pt,itemsep=2pt]
\item \textbf{Benar.} $6^3=216$.
\item \textbf{Benar.} Enumerasi solusi positif $a+b+c=7$ dengan $1\leq a,b,c\leq6$ menghasilkan 15 cara.
\item \textbf{Benar.} $P=\tfrac{15}{216}$.
\item \textbf{Salah.} $\tfrac{1}{12}=\tfrac{18}{216}\neq\tfrac{15}{216}$.
\end{enumerate}
\end{pembox}

\begin{soal}{8}
Trapesium $ABCD$: $AB=11$, $CD=7$, $BC=5$, luas $=27$.
Keliling trapesium adalah \ldots
\gambar[0.38\textwidth]{images/3/3-2.png}
\begin{pilB}\item $21$\item $24$\item $26$\item $27$\item $30$\end{pilB}
\end{soal}
\begin{pembox}
\jawaban{(C) Keliling $= 26$}
Luas: $\tfrac{1}{2}(11+7)h=9h=27\Rightarrow h=3$.
Sisi miring $BC=5$ dan tinggi $3$: selisih alas $=\sqrt{25-9}=4$.
Sisi $AD=h=3$. Keliling $=11+5+7+3=26$.
\end{pembox}

\begin{soal}{9}
Persamaan lingkaran $(x+5)^2+(y-3)^2=9$. Titik pusat dan jari-jari adalah \ldots
\begin{pil}
  \item Pusat $(5,3)$, $r=9$ \item Pusat $(-5,3)$, $r=3$
  \item Pusat $(5,3)$, $r=3$ \item Pusat $(5,-3)$, $r=9$
  \item Pusat $(-5,-3)$, $r=3$
\end{pil}
\end{soal}
\begin{pembox}
\jawaban{(B) Pusat $(-5,3)$, $r=3$}
Bentuk baku $(x-h)^2+(y-k)^2=r^2$: pusat $(-5,3)$ dan $r=\sqrt{9}=3$.
\end{pembox}

\begin{soal}{10}
$\dfrac{A}{2-\sqrt{5}}+\dfrac{B}{2+\sqrt{5}}=1+2\sqrt{5}$, $A,B\in\mathbb{R}$.
Tentukan kebenaran setiap pernyataan berikut!

\par\smallskip
{\renewcommand{\arraystretch}{1.4}%
\begin{tabular}{@{}p{0.76\linewidth}cc@{}}
\hline\textbf{Pernyataan}&\textbf{Benar}&\textbf{Salah}\\\hline
(1)~$A+B=-\dfrac{1}{2}$.&$\bigcirc$&$\bigcirc$\\
(2)~$A=-\dfrac{5}{4}$.&$\bigcirc$&$\bigcirc$\\
(3)~$B=-\dfrac{3}{4}$.&$\bigcirc$&$\bigcirc$\\
(4)~$A<B$.&$\bigcirc$&$\bigcirc$\\\hline
\end{tabular}}
\end{soal}
\begin{pembox}
\jawaban{(1)~\textbf{B}\quad(2)~\textbf{B}\quad(3)~\textbf{S}\quad(4)~\textbf{B}}
$\tfrac{1}{2-\sqrt5}=-(2+\sqrt5)$, $\tfrac{1}{2+\sqrt5}=\sqrt5-2$.
Ruas kiri $=(-2A-2B)+(-A+B)\sqrt5$. Samakan dengan $1+2\sqrt5$:
$-2(A+B)=1$ dan $-A+B=2$. Maka $A+B=-\tfrac12$, $B=A+2$, $A=-\tfrac54$, $B=\tfrac34$.
\begin{enumerate}[leftmargin=*,topsep=2pt,itemsep=2pt]
\item \textbf{Benar.} $A+B=-\tfrac12$.
\item \textbf{Benar.} $A=-\tfrac54$.
\item \textbf{Salah.} $B=\tfrac34\neq-\tfrac34$.
\item \textbf{Benar.} $-\tfrac54<\tfrac34$.
\end{enumerate}
\end{pembox}

% ===================================================================
%  SET 2  (nomor 11--20)
% ===================================================================

\begin{soal}{11}
$\begin{cases}px+qy=8\\ 3x-qy=38\end{cases}$, penyelesaian $(2,4)$. Nilai $p$ adalah \ldots
\begin{pilB}\item $40$\item $22{,}5$\item $21{,}5$\item $20$\item $8$\end{pilB}
\end{soal}
\begin{pembox}
\jawaban{(D) $p=20$}
Pers.\ 2: $6-4q=38\Rightarrow q=-8$. Pers.\ 1: $2p-32=8\Rightarrow p=20$.
\end{pembox}

\begin{soal}{12}
$a=\sqrt7+3$, $b=\sqrt7-3$. Nilai $a^3+b^3$ adalah \ldots
\begin{pilB}\item $60\sqrt7$\item $62\sqrt7$\item $64\sqrt7$\item $66\sqrt7$\item $68\sqrt7$\end{pilB}
\end{soal}
\begin{pembox}
\jawaban{(E) $68\sqrt7$}
$a+b=2\sqrt7$, $ab=7-9=-2$.
$a^3+b^3=(a+b)^3-3ab(a+b)=(2\sqrt7)^3-3(-2)(2\sqrt7)=56\sqrt7+12\sqrt7=68\sqrt7$.
\end{pembox}

\begin{soal}{13}
Solusi $\dfrac{x^2+x+3}{x^2-x-2}<0$ adalah $a<x<b$. Nilai $b-2a$ adalah \ldots
\begin{pilB}\item $1$\item $2$\item $3$\item $4$\item $5$\end{pilB}
\end{soal}
\begin{pembox}
\jawaban{(D) $b-2a=4$}
Pembilang $x^2+x+3$ selalu positif ($\Delta<0$). Pecahan negatif $\Leftrightarrow$ penyebut negatif:
$(x-2)(x+1)<0\Rightarrow-1<x<2$. Jadi $a=-1$, $b=2$, $b-2a=4$.
\end{pembox}

\begin{soal}{14}
$A=\begin{pmatrix}a&b&c\\1&-1&1\end{pmatrix}$, $B=\begin{pmatrix}1&2\\1&-1\\0&1\end{pmatrix}$.
Tentukan kebenaran setiap pernyataan berikut!

\par\smallskip
{\renewcommand{\arraystretch}{1.4}%
\begin{tabular}{@{}p{0.76\linewidth}cc@{}}
\hline\textbf{Pernyataan}&\textbf{Benar}&\textbf{Salah}\\\hline
(1)~$AB$ berukuran $2\times2$.&$\bigcirc$&$\bigcirc$\\
(2)~$\det(AB)=4(a+b)$.&$\bigcirc$&$\bigcirc$\\
(3)~Jika $\det(AB)=4$, maka $a+b=4$.&$\bigcirc$&$\bigcirc$\\
(4)~Jika $\det(AB)=4$, maka $a+b=1$.&$\bigcirc$&$\bigcirc$\\\hline
\end{tabular}}
\end{soal}
\begin{pembox}
\jawaban{(1)~\textbf{B}\quad(2)~\textbf{B}\quad(3)~\textbf{S}\quad(4)~\textbf{B}}
Baris ke-2 dari $AB$: $[0,4]$. Baris ke-1: $[a+b,\;2a-b+c]$.
$AB=\begin{pmatrix}a+b&2a-b+c\\0&4\end{pmatrix}$, $\det(AB)=4(a+b)$.
\begin{enumerate}[leftmargin=*,topsep=2pt,itemsep=2pt]
\item \textbf{Benar.} $A_{2\times3}\cdot B_{3\times2}=AB_{2\times2}$.
\item \textbf{Benar.} $\det(AB)=4(a+b)$.
\item \textbf{Salah.} $4(a+b)=4\Rightarrow a+b=1$, bukan 4.
\item \textbf{Benar.} $a+b=1$.
\end{enumerate}
\end{pembox}

\begin{soal}{15}
$ax^2+bx+b=0$ mempunyai dua akar real berbeda untuk $b<0$ atau $b>8$. Nilai $a$ adalah \ldots
\begin{pilB}\item $1$\item $2$\item $-2$\item $3$\item $-3$\end{pilB}
\end{soal}
\begin{pembox}
\jawaban{(B) $a=2$}
$\Delta=b^2-4ab=b(b-4a)>0$ berlaku untuk $b<0$ atau $b>4a$.
Dari soal batas atasnya $b>8$, sehingga $4a=8$ dan $a=2$.
\end{pembox}

\begin{soal}{16}
$f(x)=3x^2+9x+12$. Jika $\dfrac{f(x+h)-f(x)}{h}=Ax+B+Ch$,
tentukan kebenaran setiap pernyataan berikut!

\par\smallskip
{\renewcommand{\arraystretch}{1.4}%
\begin{tabular}{@{}p{0.76\linewidth}cc@{}}
\hline\textbf{Pernyataan}&\textbf{Benar}&\textbf{Salah}\\\hline
(1)~$A=6$.&$\bigcirc$&$\bigcirc$\\
(2)~$B=3$.&$\bigcirc$&$\bigcirc$\\
(3)~$C=3$.&$\bigcirc$&$\bigcirc$\\
(4)~$100A+10B+C=693$.&$\bigcirc$&$\bigcirc$\\\hline
\end{tabular}}
\end{soal}
\begin{pembox}
\jawaban{(1)~\textbf{B}\quad(2)~\textbf{S}\quad(3)~\textbf{B}\quad(4)~\textbf{B}}
$f(x+h)-f(x)=6xh+3h^2+9h$. Bagi $h$: $6x+9+3h$. Jadi $A=6$, $B=9$, $C=3$.
\begin{enumerate}[leftmargin=*,topsep=2pt,itemsep=2pt]
\item \textbf{Benar.} $A=6$.
\item \textbf{Salah.} $B=9\neq3$.
\item \textbf{Benar.} $C=3$.
\item \textbf{Benar.} $600+90+3=693$.
\end{enumerate}
\end{pembox}

\begin{soal}{17}
$f(x)=\begin{cases}x^2-x+2,&x\geq1\\x+1,&x<1\end{cases}$.
Nilai $\lim_{x\to1}f(x)$ adalah \ldots
\begin{pilB}\item $-2$\item $-1$\item Tidak ada\item $1$\item $2$\end{pilB}
\end{soal}
\begin{pembox}
\jawaban{(E) $2$}
Limit kiri: $\lim_{x\to1^-}(x+1)=2$.\quad
Limit kanan: $\lim_{x\to1^+}(x^2-x+2)=2$.
Keduanya sama, sehingga limit $=2$.
\end{pembox}

\begin{soal}{18}
$f(3)=6$, $f'(3)=4$, $g(3)=10$, $g'(3)=-6$. Nilai $\left(\dfrac{f}{g}\right)'(3)$ adalah \ldots
\begin{pilB}\item $0{,}76$\item $0{,}82$\item $0{,}91$\item $0{,}95$\item $1$\end{pilB}
\end{soal}
\begin{pembox}
\jawaban{(A) $0{,}76$}
$\left(\dfrac{f}{g}\right)'(3)=\dfrac{f'(3)g(3)-f(3)g'(3)}{[g(3)]^2}=\dfrac{4(10)-6(-6)}{100}=\dfrac{76}{100}=0{,}76$.
\end{pembox}

\begin{soal}{19}
$\displaystyle\int_1^9 f(x)\,dx=5$.
Tentukan kebenaran setiap pernyataan berikut!

\par\smallskip
{\renewcommand{\arraystretch}{1.4}%
\begin{tabular}{@{}p{0.76\linewidth}cc@{}}
\hline\textbf{Pernyataan}&\textbf{Benar}&\textbf{Salah}\\\hline
(1)~Substitusi $u=x^2$ mengubah $\int_1^3$ menjadi $\int_1^9$.&$\bigcirc$&$\bigcirc$\\
(2)~$\displaystyle\int_1^3 8x\,f(x^2)\,dx=4\int_1^9 f(u)\,du$.&$\bigcirc$&$\bigcirc$\\
(3)~Nilai $\displaystyle\int_1^3 8x\,f(x^2)\,dx=5$.&$\bigcirc$&$\bigcirc$\\
(4)~Nilai $\displaystyle\int_1^3 8x\,f(x^2)\,dx=20$.&$\bigcirc$&$\bigcirc$\\\hline
\end{tabular}}
\end{soal}
\begin{pembox}
\jawaban{(1)~\textbf{B}\quad(2)~\textbf{B}\quad(3)~\textbf{S}\quad(4)~\textbf{B}}
$u=x^2$, $du=2x\,dx$, $8x\,dx=4\,du$. Batas: $x{=}1\to u{=}1$; $x{=}3\to u{=}9$.
\begin{enumerate}[leftmargin=*,topsep=2pt,itemsep=2pt]
\item \textbf{Benar.} Batas menjadi $[1,9]$.
\item \textbf{Benar.} $\displaystyle\int_1^3 8x\,f(x^2)\,dx=4\int_1^9 f(u)\,du$.
\item \textbf{Salah.} Nilainya $4\cdot5=20$, bukan 5.
\item \textbf{Benar.} $4\cdot5=20$.
\end{enumerate}
\end{pembox}

\begin{soal}{20}
Pertidaksamaan $|x-7|\geq3$.
Tentukan kebenaran setiap pernyataan berikut!

\par\smallskip
{\renewcommand{\arraystretch}{1.4}%
\begin{tabular}{@{}p{0.76\linewidth}cc@{}}
\hline\textbf{Pernyataan}&\textbf{Benar}&\textbf{Salah}\\\hline
(1)~Bilangan yang \textbf{tidak} memenuhi $|x-7|\geq3$ memenuhi $|x-7|<3$.&$\bigcirc$&$\bigcirc$\\
(2)~$|x-7|<3$ setara dengan $4<x<10$.&$\bigcirc$&$\bigcirc$\\
(3)~Banyak bilangan bulat yang tidak memenuhi adalah $4$.&$\bigcirc$&$\bigcirc$\\
(4)~Banyak bilangan bulat yang tidak memenuhi adalah $5$.&$\bigcirc$&$\bigcirc$\\\hline
\end{tabular}}
\end{soal}
\begin{pembox}
\jawaban{(1)~\textbf{B}\quad(2)~\textbf{B}\quad(3)~\textbf{S}\quad(4)~\textbf{B}}
\begin{enumerate}[leftmargin=*,topsep=2pt,itemsep=2pt]
\item \textbf{Benar.} Negasi $\geq3$ adalah $<3$.
\item \textbf{Benar.} $|x-7|<3\Leftrightarrow4<x<10$.
\item \textbf{Salah.} Bilangan bulat dalam $(4,10)$: $\{5,6,7,8,9\}$, berjumlah 5.
\item \textbf{Benar.} Ada 5 bilangan bulat.
\end{enumerate}
\end{pembox}

% ===================================================================
%  SET 3  (nomor 21--30)
% ===================================================================

\begin{soal}{21}
Jika $9^m=4$, maka $4\cdot3^{m+1}-27^m$ adalah \ldots
\begin{pilB}\item $8$\item $12$\item $16$\item $18$\item $24$\end{pilB}
\end{soal}
\begin{pembox}
\jawaban{(C) $16$}
$3^{2m}=4\Rightarrow3^m=2$.\quad$4\cdot3^{m+1}=12\cdot2=24$.\quad$27^m=(3^m)^3=8$.
$24-8=16$.
\end{pembox}

\begin{soal}{22}
$\dfrac{1}{\log_{10}a}+\dfrac{1}{\log_{\sqrt{10}}a}+\dfrac{1}{\log_{\sqrt{\sqrt{10}}}a}+\cdots=200$.
Nilai $a$ adalah \ldots
\begin{pilB}
  \item $\tfrac{1}{100}$\item $\tfrac{1}{10}$\item $\sqrt{10}$
  \item $10^{1/100}$\item $10^{1/10}$
\end{pilB}
\end{soal}
\begin{pembox}
\jawaban{(D) $10^{1/100}$}
$\tfrac{1}{\log_b a}=\log_a b$. Basis: $10,10^{1/2},10^{1/4},\ldots$, sehingga
jumlah $=\log_a 10\!\left(1+\tfrac12+\tfrac14+\cdots\right)=2\log_a 10=200$.
$\log_a 10=100\Rightarrow a^{100}=10\Rightarrow a=10^{1/100}$.
\end{pembox}

\begin{soal}{23}
$\sin x=\dfrac{a-1{,}5}{2-0{,}5a}$.
Tentukan kebenaran setiap pernyataan berikut!

\par\smallskip
{\renewcommand{\arraystretch}{1.4}%
\begin{tabular}{@{}p{0.76\linewidth}cc@{}}
\hline\textbf{Pernyataan}&\textbf{Benar}&\textbf{Salah}\\\hline
(1)~Ruas kanan $=\dfrac{2a-3}{4-a}$.&$\bigcirc$&$\bigcirc$\\
(2)~Agar ada solusi, $-1\leq\dfrac{2a-3}{4-a}\leq1$.&$\bigcirc$&$\bigcirc$\\
(3)~Banyak bilangan bulat $a$ yang memenuhi adalah $3$.&$\bigcirc$&$\bigcirc$\\
(4)~Banyak bilangan bulat $a$ yang memenuhi adalah $4$.&$\bigcirc$&$\bigcirc$\\\hline
\end{tabular}}
\end{soal}
\begin{pembox}
\jawaban{(1)~\textbf{B}\quad(2)~\textbf{B}\quad(3)~\textbf{S}\quad(4)~\textbf{B}}
Untuk $a<4$: $-1\leq\tfrac{2a-3}{4-a}\leq1\Rightarrow-1\leq a\leq\tfrac73$.
Bilangan bulat: $-1,0,1,2$ --- ada 4.
\begin{enumerate}[leftmargin=*,topsep=2pt,itemsep=2pt]
\item \textbf{Benar.} $\tfrac{a-1{,}5}{2-0{,}5a}=\tfrac{2a-3}{4-a}$.
\item \textbf{Benar.} Syarat $\sin x\in[-1,1]$.
\item \textbf{Salah.} Ada 4 bilangan bulat, bukan 3.
\item \textbf{Benar.} $\{-1,0,1,2\}$, empat bilangan.
\end{enumerate}
\end{pembox}

\begin{soal}{24}
$\displaystyle\sum_{k=1}^{100}2k+\sum_{k=1}^{100}(3k+2)$.
Tentukan kebenaran setiap pernyataan berikut!

\par\smallskip
{\renewcommand{\arraystretch}{1.4}%
\begin{tabular}{@{}p{0.76\linewidth}cc@{}}
\hline\textbf{Pernyataan}&\textbf{Benar}&\textbf{Salah}\\\hline
(1)~$\displaystyle\sum_{k=1}^{100}k=4950$.&$\bigcirc$&$\bigcirc$\\
(2)~$\displaystyle\sum_{k=1}^{100}2k=10100$.&$\bigcirc$&$\bigcirc$\\
(3)~$\displaystyle\sum_{k=1}^{100}(3k+2)=15350$.&$\bigcirc$&$\bigcirc$\\
(4)~Total $=25450$.&$\bigcirc$&$\bigcirc$\\\hline
\end{tabular}}
\end{soal}
\begin{pembox}
\jawaban{(1)~\textbf{S}\quad(2)~\textbf{B}\quad(3)~\textbf{B}\quad(4)~\textbf{B}}
$\sum_{k=1}^{100}k=\tfrac{100\cdot101}{2}=5050$.
\begin{enumerate}[leftmargin=*,topsep=2pt,itemsep=2pt]
\item \textbf{Salah.} $\sum k=5050\neq4950$.
\item \textbf{Benar.} $2\times5050=10100$.
\item \textbf{Benar.} $3(5050)+2(100)=15150+200=15350$.
\item \textbf{Benar.} $10100+15350=25450$.
\end{enumerate}
\end{pembox}

\begin{soal}{25}
$(x+1)(x^2+2x-7)\geq x^2-1$.
Tentukan kebenaran setiap pernyataan berikut!

\par\smallskip
{\renewcommand{\arraystretch}{1.4}%
\begin{tabular}{@{}p{0.76\linewidth}cc@{}}
\hline\textbf{Pernyataan}&\textbf{Benar}&\textbf{Salah}\\\hline
(1)~Dapat difaktorkan menjadi $(x-2)(x+1)(x+3)\geq0$.&$\bigcirc$&$\bigcirc$\\
(2)~$x=-3$ memenuhi.&$\bigcirc$&$\bigcirc$\\
(3)~Banyak bilangan bulat negatif yang memenuhi adalah $3$.&$\bigcirc$&$\bigcirc$\\
(4)~$x=-4$ memenuhi.&$\bigcirc$&$\bigcirc$\\\hline
\end{tabular}}
\end{soal}
\begin{pembox}
\jawaban{(1)~\textbf{B}\quad(2)~\textbf{B}\quad(3)~\textbf{B}\quad(4)~\textbf{S}}
$(x+1)(x^2+2x-7)-(x^2-1)=x^3+2x^2-5x-6=(x-2)(x+1)(x+3)\geq0$.
Tanda positif pada $[-3,-1]\cup[2,\infty)$.
\begin{enumerate}[leftmargin=*,topsep=2pt,itemsep=2pt]
\item \textbf{Benar.} Faktorisasi benar.
\item \textbf{Benar.} $x=-3\in[-3,-1]$.
\item \textbf{Benar.} Bilangan bulat negatif: $-3,-2,-1$.
\item \textbf{Salah.} $(-6)(-3)(-1)=-18<0$.
\end{enumerate}
\end{pembox}

\begin{soal}{26}
Penjualan 10 hari terakhir: $8,7,8,6,8,8,9,8,10,8$. Rata-rata bulan Juni $=10$.
Rata-rata 20 hari pertama adalah \ldots
\begin{pilB}\item $8$\item $9$\item $10$\item $11$\item $13$\end{pilB}
\end{soal}
\begin{pembox}
\jawaban{(D) $11$}
Total 30 hari: $300$. Jumlah 10 hari terakhir: $80$. Jumlah 20 hari pertama: $220$.
Rata-rata: $220/20=11$.
\end{pembox}

\begin{soal}{27}
$x^7-10x^6+24x^5=0$, tiga akar real $r_1<r_2<r_3$.
Tentukan kebenaran setiap pernyataan berikut!

\par\smallskip
{\renewcommand{\arraystretch}{1.4}%
\begin{tabular}{@{}p{0.76\linewidth}cc@{}}
\hline\textbf{Pernyataan}&\textbf{Benar}&\textbf{Salah}\\\hline
(1)~$x^5(x-4)(x-6)=0$.&$\bigcirc$&$\bigcirc$\\
(2)~Tiga akar berbeda: $0,4,6$.&$\bigcirc$&$\bigcirc$\\
(3)~$r_1+r_3=6$.&$\bigcirc$&$\bigcirc$\\
(4)~$r_1\cdot r_3=24$.&$\bigcirc$&$\bigcirc$\\\hline
\end{tabular}}
\end{soal}
\begin{pembox}
\jawaban{(1)~\textbf{B}\quad(2)~\textbf{B}\quad(3)~\textbf{B}\quad(4)~\textbf{S}}
\begin{enumerate}[leftmargin=*,topsep=2pt,itemsep=2pt]
\item \textbf{Benar.} $x^5(x^2-10x+24)=x^5(x-4)(x-6)$.
\item \textbf{Benar.} $r_1=0,r_2=4,r_3=6$.
\item \textbf{Benar.} $0+6=6$.
\item \textbf{Salah.} $0\times6=0\neq24$.
\end{enumerate}
\end{pembox}

\begin{soal}{28}
Garis singgung $y=3x^2+C$ di suatu titik adalah $y=24x$. Nilai $C$ adalah \ldots
\begin{pilB}\item $50$\item $46$\item $47$\item $48$\item $49$\end{pilB}
\end{soal}
\begin{pembox}
\jawaban{(D) $C=48$}
$y'=6x=24\Rightarrow x=4$. Pada garis: $y=96$. Pada kurva: $48+C=96\Rightarrow C=48$.
\end{pembox}

\begin{soal}{29}
Transformasi parabola $y=x^2$.
Tentukan kebenaran setiap pernyataan berikut!

\par\smallskip
{\renewcommand{\arraystretch}{1.4}%
\begin{tabular}{@{}p{0.76\linewidth}cc@{}}
\hline\textbf{Pernyataan}&\textbf{Benar}&\textbf{Salah}\\\hline
(1)~$y=(x-h)^2$ adalah pergeseran $y=x^2$ sejauh $h$ ke kanan ($h>0$).&$\bigcirc$&$\bigcirc$\\
(2)~$y=(x-1)^2$ diperoleh dengan menggeser $y=x^2$ ke kanan 1 satuan.&$\bigcirc$&$\bigcirc$\\
(3)~$y=(x-1)^2$ diperoleh dengan menggeser $y=x^2$ ke kiri 1 satuan.&$\bigcirc$&$\bigcirc$\\
(4)~Titik minimum $y=(x-1)^2$ berada di $(1,0)$.&$\bigcirc$&$\bigcirc$\\\hline
\end{tabular}}
\end{soal}
\begin{pembox}
\jawaban{(1)~\textbf{B}\quad(2)~\textbf{B}\quad(3)~\textbf{S}\quad(4)~\textbf{B}}
\begin{enumerate}[leftmargin=*,topsep=2pt,itemsep=2pt]
\item \textbf{Benar.} Definisi pergeseran horizontal.
\item \textbf{Benar.} $h=1$, ke kanan 1 satuan.
\item \textbf{Salah.} Ke kanan, bukan ke kiri.
\item \textbf{Benar.} Minimum saat $x=1$: $y=0$.
\end{enumerate}
\end{pembox}

\begin{soal}{30}
Bilangan bulat terdekat dengan $5\sqrt{2}$ adalah \ldots\hfill
Jawaban:\enspace\underline{\hspace{3cm}}
\end{soal}
\begin{pembox}
\jawaban{$7$}
$5\sqrt{2}\approx5\times1{,}414=7{,}07$. Bilangan bulat terdekat: $7$.
\end{pembox}

% ===================================================================
%  SET 4  (nomor 31--40)
% ===================================================================

\begin{soal}{31}
$f(x)=\sqrt{1+x}$. Perhatikan $\displaystyle\lim_{h\to0}\dfrac{f(3+2h^2)-f(3-3h^2)}{h^2}$.
Tentukan kebenaran setiap pernyataan berikut!

\par\smallskip
{\renewcommand{\arraystretch}{1.4}%
\begin{tabular}{@{}p{0.76\linewidth}cc@{}}
\hline\textbf{Pernyataan}&\textbf{Benar}&\textbf{Salah}\\\hline
(1)~$f'(x)=\dfrac{1}{2\sqrt{1+x}}$.&$\bigcirc$&$\bigcirc$\\
(2)~$f'(3)=\dfrac{1}{4}$.&$\bigcirc$&$\bigcirc$\\
(3)~Nilai limit adalah $\dfrac{9}{8}$.&$\bigcirc$&$\bigcirc$\\
(4)~Nilai limit adalah $\dfrac{5}{4}$.&$\bigcirc$&$\bigcirc$\\\hline
\end{tabular}}
\end{soal}
\begin{pembox}
\jawaban{(1)~\textbf{B}\quad(2)~\textbf{B}\quad(3)~\textbf{S}\quad(4)~\textbf{B}}
$f(3+2h^2)-f(3-3h^2)\approx f'(3)\cdot5h^2$.
\begin{enumerate}[leftmargin=*,topsep=2pt,itemsep=2pt]
\item \textbf{Benar.} Turunan benar.
\item \textbf{Benar.} $f'(3)=\tfrac{1}{2\sqrt4}=\tfrac14$.
\item \textbf{Salah.} Limitnya $\tfrac14\times5=\tfrac54\neq\tfrac98$.
\item \textbf{Benar.} Nilai limit $=\dfrac{5}{4}$.
\end{enumerate}
\end{pembox}

\begin{soal}{32}
$\cos(3x+15^\circ)=\sin(x+25^\circ)$, $0^\circ<x<90^\circ$.
Tentukan kebenaran setiap pernyataan berikut!

\par\smallskip
{\renewcommand{\arraystretch}{1.4}%
\begin{tabular}{@{}p{0.76\linewidth}cc@{}}
\hline\textbf{Pernyataan}&\textbf{Benar}&\textbf{Salah}\\\hline
(1)~$\sin(x+25^\circ)=\cos(65^\circ-x)$.&$\bigcirc$&$\bigcirc$\\
(2)~Dari persamaan: $4x=50^\circ$.&$\bigcirc$&$\bigcirc$\\
(3)~$x=15^\circ$.&$\bigcirc$&$\bigcirc$\\
(4)~$x=12{,}5^\circ$.&$\bigcirc$&$\bigcirc$\\\hline
\end{tabular}}
\end{soal}
\begin{pembox}
\jawaban{(1)~\textbf{B}\quad(2)~\textbf{B}\quad(3)~\textbf{S}\quad(4)~\textbf{B}}
\begin{enumerate}[leftmargin=*,topsep=2pt,itemsep=2pt]
\item \textbf{Benar.} $\sin\theta=\cos(90^\circ-\theta)$.
\item \textbf{Benar.} $3x+15^\circ=65^\circ-x\Rightarrow4x=50^\circ$.
\item \textbf{Salah.} $x=12{,}5^\circ\neq15^\circ$.
\item \textbf{Benar.} $x=\tfrac{50^\circ}{4}=12{,}5^\circ$.
\end{enumerate}
\end{pembox}

\begin{soal}{33}
$f(x)=\sqrt{\dfrac{x^2-5x}{1-x}}$.
Tentukan kebenaran setiap pernyataan berikut!

\par\smallskip
{\renewcommand{\arraystretch}{1.4}%
\begin{tabular}{@{}p{0.76\linewidth}cc@{}}
\hline\textbf{Pernyataan}&\textbf{Benar}&\textbf{Salah}\\\hline
(1)~Diperlukan $\dfrac{x(x-5)}{1-x}\geq0$.&$\bigcirc$&$\bigcirc$\\
(2)~$1<x\leq5$ adalah bagian daerah asal.&$\bigcirc$&$\bigcirc$\\
(3)~$0<x<1$ adalah bagian daerah asal.&$\bigcirc$&$\bigcirc$\\
(4)~Daerah asal: $\{x\mid x\leq0\text{ atau }1<x\leq5\}$.&$\bigcirc$&$\bigcirc$\\\hline
\end{tabular}}
\end{soal}
\begin{pembox}
\jawaban{(1)~\textbf{B}\quad(2)~\textbf{B}\quad(3)~\textbf{S}\quad(4)~\textbf{B}}
Uji tanda $\tfrac{x(x-5)}{1-x}$ pada titik kritis $0,1,5$:
$x<0$: $\frac{(-)(-)}{(+)}>0$\,$\checkmark$;\; $0<x<1$: $\frac{(+)(-)}{(+)}<0$\,$\times$;\;
$1<x<5$: $\frac{(+)(-)}{(-)}>0$\,$\checkmark$;\; $x>5$: $\frac{(+)(+)}{(-)}<0$\,$\times$.
\begin{enumerate}[leftmargin=*,topsep=2pt,itemsep=2pt]
\item \textbf{Benar.} Kondisi tersebut diperlukan.
\item \textbf{Benar.} $1<x\leq5$ memenuhi.
\item \textbf{Salah.} $0<x<1$ memberikan nilai negatif.
\item \textbf{Benar.} Daerah asal: $x\leq0$ atau $1<x\leq5$.
\end{enumerate}
\end{pembox}

\begin{soal}{34}
$\displaystyle\int_0^1\dfrac{x}{1+x}\,dx=a$.
Nilai $\displaystyle\int_0^1\dfrac{1}{1+x}\,dx$ adalah \ldots\hfill
Jawaban:\enspace\underline{\hspace{3cm}}
\end{soal}
\begin{pembox}
\jawaban{$1-a$}
$\dfrac{x}{1+x}+\dfrac{1}{1+x}=1$, sehingga $a+\displaystyle\int_0^1\dfrac{1}{1+x}\,dx=1$.
Jawaban: $1-a$.
\end{pembox}

\begin{soal}{35}
$f(x)\geq0$ untuk semua $x$, $f(1)=0$, $f(2)=2$. Nilai $f(0)+f(4)$ adalah \ldots\hfill
Jawaban:\enspace\underline{\hspace{3cm}}
\end{soal}
\begin{pembox}
\jawaban{$20$}
$f(1)=0$ dan $f\geq0\Rightarrow x=1$ akar kembar: $f(x)=k(x-1)^2$.
$f(2)=k=2$. Maka $f(0)=2$ dan $f(4)=2\cdot9=18$. Jumlah $=20$.
\end{pembox}

\begin{soal}{36}
$A,B$ pada lingkaran $O(0,0)$, $\alpha=\angle AOB$, $A(a,b)$ kuadran I, $B=(25,0)$,
$\sin\alpha=\tfrac{4}{5}$.
\begin{enumerate}[label=(\alph*),leftmargin=*,topsep=3pt,itemsep=3pt]
  \item Nilai $b$.\hfill Jawaban:\enspace\underline{\hspace{3.5cm}}
  \item Nilai $5\cos\alpha$.\hfill Jawaban:\enspace\underline{\hspace{3.5cm}}
\end{enumerate}
\end{soal}
\begin{pembox}
\jawaban{(a) $b=20$\quad(b) $5\cos\alpha=3$}
$r=OB=25$. $A=(25\cos\alpha,25\sin\alpha)$.
\textbf{(a)} $b=25\sin\alpha=25\cdot\tfrac45=20$.
\textbf{(b)} $\cos\alpha=\tfrac35$, $5\cos\alpha=3$.
\end{pembox}

\begin{soal}{37}
$U_6=\log_8 96$, $U_2=\log_8 6$, $U_{12}-U_4=\log_8(c)$. Nilai $c$ adalah \ldots\hfill
Jawaban:\enspace\underline{\hspace{3cm}}
\end{soal}
\begin{pembox}
\jawaban{$c=256$}
$4d=U_6-U_2=\log_8 16$.\quad$d=\tfrac14\log_8 16$.
$U_{12}-U_4=8d=2\log_8 16=\log_8(16^2)=\log_8 256$.\quad$c=256$.
\end{pembox}

\begin{soal}{38}
$f(-x)=f(x)$, $\int_0^5 f=7$, $\int_5^6 f=8$.
Nilai $\displaystyle\int_{-5}^6 f(x)\,dx$ adalah \ldots\hfill
Jawaban:\enspace\underline{\hspace{3cm}}
\end{soal}
\begin{pembox}
\jawaban{$22$}
Fungsi genap: $\int_{-5}^0 f=\int_0^5 f=7$.
$\int_{-5}^6 f=7+7+8=22$.
\end{pembox}

\begin{soal}{39}
$\displaystyle\lim_{x\to0}\dfrac{f(x)}{3x}=5$. Nilai
$\displaystyle\lim_{x\to0}\!\left[\dfrac{2f(x)}{x^2}-\dfrac{f(x)\cos2x}{x(1+x)}-\dfrac{2f(x)}{x^2(1+x)}\right]$
adalah \ldots\hfill Jawaban:\enspace\underline{\hspace{2.5cm}}
\end{soal}
\begin{pembox}
\jawaban{$15$}
$\lim\tfrac{f(x)}{x}=15$. Ekspresi $=\dfrac{f(x)}{x}\cdot\dfrac{2(1+x)-x\cos2x-2}{x(1+x)}
=\dfrac{f(x)}{x}\cdot\dfrac{2-\cos2x}{1+x}$.
Limit: $15\cdot\dfrac{2-1}{1}=15$.
\end{pembox}

\begin{soal}{40}
$l$: $y=x$; $k$: $y=-\tfrac53x$. $A$ pada $k$, $B$ proyeksi $A$ ke $l$, $|OB|=1$.
Jarak $AB$ adalah \ldots\hfill Jawaban:\enspace\underline{\hspace{3cm}}
\end{soal}
\begin{pembox}
\jawaban{$4$}
Vektor satuan pada $l$: $\mathbf{u}=\tfrac{(1,1)}{\sqrt2}$; tegak lurus: $\mathbf{n}=\tfrac{(1,-1)}{\sqrt2}$.
$B=\mathbf{u}$, $A=B+s\mathbf{n}$.
Syarat $A$ pada $k$: $\tfrac{1-s}{\sqrt2}=-\tfrac53\cdot\tfrac{1+s}{\sqrt2}\Rightarrow3-3s=-5-5s\Rightarrow s=-4$.
Jarak $AB=|s|=4$.
\end{pembox}

% ===================================================================
%  SET 5  (nomor 41--50)
% ===================================================================

\begin{soal}{41}
Truk $A,B$ memuat 4 ton; truk $C,D$ memuat 6 ton. Truk $E$ memuat
1 ton lebih dari rata-rata kelima truk. Muatan $A+E$ adalah \ldots
\begin{pilB}\item $8{,}17$\item $9$\item $10$\item $10{,}25$\item $10{,}5$\end{pilB}
\end{soal}
\begin{pembox}
\jawaban{(D) $10{,}25$ ton}
Misalkan muatan $E=e$. Rata-rata $=\dfrac{4+4+6+6+e}{5}=\dfrac{20+e}{5}$.
$e=\dfrac{20+e}{5}+1\Rightarrow5e=25+e\Rightarrow e=6{,}25$.
$A+E=4+6{,}25=10{,}25$.
\end{pembox}

\begin{soal}{42}
Satu dadu dilempar 3 kali. Peluang mata 6 muncul sedikitnya sekali adalah \ldots
\begin{pilB}
  \item $\tfrac{1}{216}$\item $\tfrac{3}{216}$\item $\tfrac{12}{216}$
  \item $\tfrac{18}{216}$\item $\tfrac{91}{216}$
\end{pilB}
\end{soal}
\begin{pembox}
\jawaban{(E) $\dfrac{91}{216}$}
Komplemen: $P(\text{tidak muncul 6 sama sekali})=\left(\tfrac56\right)^3=\tfrac{125}{216}$.
$P(\text{sedikitnya sekali})=1-\tfrac{125}{216}=\tfrac{91}{216}$.
\end{pembox}

\begin{soal}{43}
$x^3+3x^2+9x+3$ membagi habis $x^4+4x^3+2ax^2+4bx+c$.
Tentukan kebenaran setiap pernyataan berikut!

\par\smallskip
{\renewcommand{\arraystretch}{1.4}%
\begin{tabular}{@{}p{0.76\linewidth}cc@{}}
\hline\textbf{Pernyataan}&\textbf{Benar}&\textbf{Salah}\\\hline
(1)~Hasil bagi berbentuk $x+k$ untuk suatu $k$.&$\bigcirc$&$\bigcirc$\\
(2)~Nilai $k=1$.&$\bigcirc$&$\bigcirc$\\
(3)~Nilai $a=6$.&$\bigcirc$&$\bigcirc$\\
(4)~Nilai $b=6$.&$\bigcirc$&$\bigcirc$\\\hline
\end{tabular}}
\end{soal}
\begin{pembox}
\jawaban{(1)~\textbf{B}\quad(2)~\textbf{B}\quad(3)~\textbf{B}\quad(4)~\textbf{S}}
Pembagi berderajat 3, yang dibagi berderajat 4, sehingga hasil bagi berderajat 1: $x+k$.
$(x^3+3x^2+9x+3)(x+k)=x^4+(k+3)x^3+(3k+9)x^2+(9k+3)x+3k$.
Samakan koefisien:
$k+3=4\Rightarrow k=1$;\quad$3k+9=2a\Rightarrow a=6$;\quad$9k+3=4b\Rightarrow b=3$.
\begin{enumerate}[leftmargin=*,topsep=2pt,itemsep=2pt]
\item \textbf{Benar.} Hasil bagi berbentuk $x+k$.
\item \textbf{Benar.} $k=1$.
\item \textbf{Benar.} $a=6$.
\item \textbf{Salah.} $b=3\neq6$.
\end{enumerate}
\end{pembox}

\begin{soal}{44}
$y=2x^2+1$ dicerminkan terhadap sumbu $x$, digeser kanan 2 satuan $\Rightarrow y=ax^2+bx+c$.
Tentukan kebenaran setiap pernyataan berikut!

\par\smallskip
{\renewcommand{\arraystretch}{1.4}%
\begin{tabular}{@{}p{0.76\linewidth}cc@{}}
\hline\textbf{Pernyataan}&\textbf{Benar}&\textbf{Salah}\\\hline
(1)~Pencerminan terhadap sumbu $x$ menghasilkan $y=-2x^2-1$.&$\bigcirc$&$\bigcirc$\\
(2)~Pergeseran kanan 2 satuan mengganti $x$ dengan $x-2$.&$\bigcirc$&$\bigcirc$\\
(3)~Persamaan akhirnya $y=-2x^2+8x-7$.&$\bigcirc$&$\bigcirc$\\
(4)~Nilai $a+b+c=-3$.&$\bigcirc$&$\bigcirc$\\\hline
\end{tabular}}
\end{soal}
\begin{pembox}
\jawaban{(1)~\textbf{B}\quad(2)~\textbf{B}\quad(3)~\textbf{S}\quad(4)~\textbf{B}}
Cermin $x$-axis: $y=-2x^2-1$. Geser kanan 2: ganti $x\to x-2$:
$y=-2(x-2)^2-1=-2x^2+8x-8-1=-2x^2+8x-9$.
Jadi $a=-2$, $b=8$, $c=-9$.
\begin{enumerate}[leftmargin=*,topsep=2pt,itemsep=2pt]
\item \textbf{Benar.} $y=-2x^2-1$.
\item \textbf{Benar.} $x\to x-2$.
\item \textbf{Salah.} Persamaan akhir adalah $-2x^2+8x-9$, bukan $-7$.
\item \textbf{Benar.} $-2+8-9=-3$.
\end{enumerate}
\end{pembox}

\begin{soal}{45}
Nilai $\sqrt[3]{9+4\sqrt{5}}+\sqrt[3]{9-4\sqrt{5}}$ adalah \ldots\hfill
Jawaban:\enspace\underline{\hspace{3cm}}
\end{soal}
\begin{pembox}
\jawaban{$3$}
Misalkan $s=\sqrt[3]{9+4\sqrt5}+\sqrt[3]{9-4\sqrt5}$.
Perkalian: $\sqrt[3]{(9+4\sqrt5)(9-4\sqrt5)}=\sqrt[3]{81-80}=1$.
$s^3=(9+4\sqrt5)+(9-4\sqrt5)+3\cdot1\cdot s=18+3s$.
$s^3-3s-18=0$; coba $s=3$: $27-9-18=0$ ✓. Jadi $s=3$.
\end{pembox}

\begin{soal}{46}
Nilai $\displaystyle\lim_{x\to5}\dfrac{\sin(30-6x)}{x^2-25}$ adalah \ldots
\begin{pilB}\item $-0{,}2$\item $-0{,}4$\item $-0{,}6$\item $-0{,}8$\item $-1$\end{pilB}
\end{soal}
\begin{pembox}
\jawaban{(C) $-0{,}6$}
Saat $x\to5$: $30-6x\to0$, sehingga $\sin(30-6x)\approx30-6x=-6(x-5)$.
$x^2-25=(x-5)(x+5)$.
$\displaystyle\lim_{x\to5}\dfrac{-6(x-5)}{(x-5)(x+5)}=\dfrac{-6}{10}=-0{,}6$.
\end{pembox}

\begin{soal}{47}
Parabola melalui $(-4,0)$ dan $(2,0)$, titik minimum pada $y=23x+20$.
Tentukan kebenaran setiap pernyataan berikut!

\par\smallskip
{\renewcommand{\arraystretch}{1.4}%
\begin{tabular}{@{}p{0.76\linewidth}cc@{}}
\hline\textbf{Pernyataan}&\textbf{Benar}&\textbf{Salah}\\\hline
(1)~Sumbu simetri berada di $x=-1$.&$\bigcirc$&$\bigcirc$\\
(2)~Ordinat titik minimum adalah $-3$.&$\bigcirc$&$\bigcirc$\\
(3)~Nilai $k=\dfrac{1}{3}$ dalam $y=k(x+4)(x-2)$.&$\bigcirc$&$\bigcirc$\\
(4)~Persamaan parabolanya $y=\dfrac{1}{3}(x-4)(x+2)$.&$\bigcirc$&$\bigcirc$\\\hline
\end{tabular}}
\end{soal}
\begin{pembox}
\jawaban{(1)~\textbf{B}\quad(2)~\textbf{B}\quad(3)~\textbf{B}\quad(4)~\textbf{S}}
Bentuk: $y=k(x+4)(x-2)$. Sumbu simetri: $x=\tfrac{-4+2}{2}=-1$.
Pada garis $y=23x+20$ saat $x=-1$: $y=-23+20=-3$ (ordinat minimum).
Pada parabola saat $x=-1$: $k(3)(-3)=-9k=-3\Rightarrow k=\tfrac13$.
\begin{enumerate}[leftmargin=*,topsep=2pt,itemsep=2pt]
\item \textbf{Benar.} Sumbu simetri di $x=-1$.
\item \textbf{Benar.} Ordinat minimum $=-3$.
\item \textbf{Benar.} $k=\tfrac13$.
\item \textbf{Salah.} Persamaan yang benar $y=\tfrac13(x+4)(x-2)$, bukan $(x-4)(x+2)$.
\end{enumerate}
\end{pembox}

\begin{soal}{48}
Rata-rata $5,4,c+11,c-1,5$ adalah $10$. Nilai $c$ adalah \ldots
\begin{pilB}\item $18$\item $17$\item $16$\item $15$\item $13$\end{pilB}
\end{soal}
\begin{pembox}
\jawaban{(E) $c=13$}
Jumlah $=5\times10=50$. $5+4+(c+11)+(c-1)+5=2c+24=50\Rightarrow c=13$.
\end{pembox}

\begin{soal}{49}
Data terurut: $4,\,6-a,\,6,\,8,\,8+a,\,10$ dengan $0\leq a\leq2$,
$P=s^2$ (varians), $Q=5$.
Tentukan kebenaran setiap pernyataan berikut!

\par\smallskip
{\renewcommand{\arraystretch}{1.4}%
\begin{tabular}{@{}p{0.76\linewidth}cc@{}}
\hline\textbf{Pernyataan}&\textbf{Benar}&\textbf{Salah}\\\hline
(1)~Rata-rata data adalah $7$ untuk semua $a$.&$\bigcirc$&$\bigcirc$\\
(2)~$s^2=\dfrac{11+2a+a^2}{3}$.&$\bigcirc$&$\bigcirc$\\
(3)~$P<Q$ untuk semua $0\leq a\leq2$.&$\bigcirc$&$\bigcirc$\\
(4)~Hubungan $P$ dan $Q$ tidak dapat ditentukan secara pasti.&$\bigcirc$&$\bigcirc$\\\hline
\end{tabular}}
\end{soal}
\begin{pembox}
\jawaban{(1)~\textbf{B}\quad(2)~\textbf{B}\quad(3)~\textbf{S}\quad(4)~\textbf{B}}
Jumlah data $=42$; rata-rata $=7$ (konstan). Deviasi dari 7: $-3,-(1+a),-1,1,(1+a),3$.
$s^2=\tfrac{9+(1+a)^2+1+1+(1+a)^2+9}{6}=\tfrac{10+(1+a)^2}{3}=\tfrac{11+2a+a^2}{3}$.
\begin{enumerate}[leftmargin=*,topsep=2pt,itemsep=2pt]
\item \textbf{Benar.} Rata-rata $=7$ untuk semua $a$.
\item \textbf{Benar.} $s^2=\tfrac{11+2a+a^2}{3}$.
\item \textbf{Salah.} Untuk $a=0$: $s^2=\tfrac{11}{3}\approx3{,}67<5$. Untuk $a=2$: $s^2=\tfrac{19}{3}\approx6{,}33>5$.
\item \textbf{Benar.} Bergantung $a$, hubungan $P$ dan $Q$ tidak pasti.
\end{enumerate}
\end{pembox}

\begin{soal}{50}
Data $D_0$ terdiri atas tiga bilangan positif dengan median 20.
Empat bilangan positif $a,b,c,d$ ditambahkan membentuk $D$ berisi 7
bilangan. Jika $a+b=40$ dan $c+d=40$, tentukan banyaknya
\textbf{maksimum} bilangan dalam $D$ yang nilainya kurang dari 20.\\[0.4em]
Jawaban:\enspace\underline{\hspace{4cm}}
\end{soal}
\begin{pembox}
\jawaban{$3$}
Dari $D_0$: median $=20$, sehingga bilangan terkecilnya $\leq20$.
Paling banyak \textbf{1} anggota $D_0$ yang nilainya $<20$.

Dari $\{a,b\}$ dengan $a+b=40$ (positif): jika $a<20$ maka $b=40-a>20$.
Paling banyak \textbf{1} dari $\{a,b\}$ yang $<20$. Demikian pula $\{c,d\}$.

Maksimum bilangan dalam $D$ yang $<20$: $1+1+1=\mathbf{3}$.

Karena median 7 bilangan adalah bilangan ke-4 dan paling banyak 3 bilangan $<20$,
bilangan ke-4 terkecil $\geq20$, sehingga median $D\geq20$.
\end{pembox}

% ===================================================================
%  SET 6  (nomor 51--60)
% ===================================================================

\begin{soal}{51}
Nilai maksimum $\dfrac{m}{15\sin x-8\cos x+25}=2$. Nilai $m$ adalah \ldots
\begin{pilB}\item $4$\item $16$\item $36$\item $64$\item $84$\end{pilB}
\end{soal}
\begin{pembox}
\jawaban{(B) $m=16$}
$15\sin x-8\cos x$ memiliki nilai minimum $-\sqrt{15^2+8^2}=-17$.
Penyebut minimum $=25-17=8$.
Untuk $m>0$: nilai maksimum pecahan $=\tfrac{m}{8}=2\Rightarrow m=16$.
\end{pembox}

\begin{soal}{52}
$f(x^2+3x+1)=\log_2(2x^3-x^2+7)$, $x\geq0$. Nilai $f(5)$ adalah \ldots
\begin{pilB}\item $1$\item $2$\item $3$\item $4$\item $5$\end{pilB}
\end{soal}
\begin{pembox}
\jawaban{(C) $f(5)=3$}
Cari $x\geq0$ dengan $x^2+3x+1=5$: $x^2+3x-4=0\Rightarrow(x-1)(x+4)=0\Rightarrow x=1$.
$f(5)=\log_2(2-1+7)=\log_2 8=3$.
\end{pembox}

\begin{soal}{53}
$\displaystyle\lim_{x\to\infty}\!\left(\sqrt{4x^2+8x}-\sqrt{x^2+1}-\sqrt{x^2+x}\right)$.
Tentukan kebenaran setiap pernyataan berikut!

\par\smallskip
{\renewcommand{\arraystretch}{1.4}%
\begin{tabular}{@{}p{0.76\linewidth}cc@{}}
\hline\textbf{Pernyataan}&\textbf{Benar}&\textbf{Salah}\\\hline
(1)~Untuk $x\to\infty$: $\sqrt{4x^2+8x}=2x+2+o(1)$.&$\bigcirc$&$\bigcirc$\\
(2)~Untuk $x\to\infty$: $\sqrt{x^2+x}=x+\tfrac12+o(1)$.&$\bigcirc$&$\bigcirc$\\
(3)~Nilai limit adalah $2$.&$\bigcirc$&$\bigcirc$\\
(4)~Nilai limit adalah $\dfrac{3}{2}$.&$\bigcirc$&$\bigcirc$\\\hline
\end{tabular}}
\end{soal}
\begin{pembox}
\jawaban{(1)~\textbf{B}\quad(2)~\textbf{B}\quad(3)~\textbf{S}\quad(4)~\textbf{B}}
Untuk $x\to\infty$: $\sqrt{a^2x^2+bx}\approx ax+\tfrac{b}{2a}$.
\begin{enumerate}[leftmargin=*,topsep=2pt,itemsep=2pt]
\item \textbf{Benar.} $\sqrt{4x^2+8x}=2\sqrt{x^2+2x}\approx2(x+1)=2x+2$.
\item \textbf{Benar.} $\sqrt{x^2+x}\approx x+\tfrac12$.
\item \textbf{Salah.} Limit $=(2x+2)-x-(x+\tfrac12)=\tfrac32\neq2$.
\item \textbf{Benar.} Nilai limit $=\dfrac{3}{2}$.
\end{enumerate}
\end{pembox}

\begin{soal}{54}
$y=\tfrac13x^3-\tfrac32x^2+2x$ mempunyai garis singgung mendatar di $P$ dan $Q$.
Tentukan kebenaran setiap pernyataan berikut!

\par\smallskip
{\renewcommand{\arraystretch}{1.4}%
\begin{tabular}{@{}p{0.76\linewidth}cc@{}}
\hline\textbf{Pernyataan}&\textbf{Benar}&\textbf{Salah}\\\hline
(1)~$y'=x^2-3x+2=0$ di $x=1$ dan $x=2$.&$\bigcirc$&$\bigcirc$\\
(2)~Ordinat di $x=1$ adalah $\tfrac{5}{6}$ dan di $x=2$ adalah $\tfrac{2}{3}$.&$\bigcirc$&$\bigcirc$\\
(3)~Jumlah ordinat $P$ dan $Q$ adalah $\dfrac{5}{3}$.&$\bigcirc$&$\bigcirc$\\
(4)~Jumlah ordinat $P$ dan $Q$ adalah $\dfrac{3}{2}$.&$\bigcirc$&$\bigcirc$\\\hline
\end{tabular}}
\end{soal}
\begin{pembox}
\jawaban{(1)~\textbf{B}\quad(2)~\textbf{B}\quad(3)~\textbf{S}\quad(4)~\textbf{B}}
$y'=x^2-3x+2=(x-1)(x-2)=0$ di $x=1,2$.
$y(1)=\tfrac13-\tfrac32+2=\tfrac{2-9+12}{6}=\tfrac56$.\quad
$y(2)=\tfrac83-6+4=\tfrac83-2=\tfrac23$.
Jumlah $=\tfrac56+\tfrac23=\tfrac56+\tfrac46=\tfrac96=\tfrac32$.
\begin{enumerate}[leftmargin=*,topsep=2pt,itemsep=2pt]
\item \textbf{Benar.} $x=1$ dan $x=2$.
\item \textbf{Benar.} Ordinat benar.
\item \textbf{Salah.} Jumlah $=\tfrac32\neq\tfrac53$.
\item \textbf{Benar.} Jumlah $=\dfrac32$.
\end{enumerate}
\end{pembox}

\begin{soal}{55}
Lingkaran pusat $(a,b)$, $a,b>3$, menyinggung $3x+4y=12$, jari-jari $12$.
Tentukan kebenaran setiap pernyataan berikut!

\par\smallskip
{\renewcommand{\arraystretch}{1.4}%
\begin{tabular}{@{}p{0.76\linewidth}cc@{}}
\hline\textbf{Pernyataan}&\textbf{Benar}&\textbf{Salah}\\\hline
(1)~Jarak pusat ke garis $=\dfrac{|3a+4b-12|}{5}$.&$\bigcirc$&$\bigcirc$\\
(2)~Karena $a,b>3$: $3a+4b-12>0$.&$\bigcirc$&$\bigcirc$\\
(3)~$3a+4b=60$.&$\bigcirc$&$\bigcirc$\\
(4)~$3a+4b=72$.&$\bigcirc$&$\bigcirc$\\\hline
\end{tabular}}
\end{soal}
\begin{pembox}
\jawaban{(1)~\textbf{B}\quad(2)~\textbf{B}\quad(3)~\textbf{S}\quad(4)~\textbf{B}}
Jarak $=r$: $\dfrac{|3a+4b-12|}{5}=12$. Karena $a,b>3$: $3a+4b>9+12=21>12$, sehingga nilai mutlak positif.
$3a+4b-12=60\Rightarrow3a+4b=72$.
\begin{enumerate}[leftmargin=*,topsep=2pt,itemsep=2pt]
\item \textbf{Benar.} Rumus jarak titik ke garis.
\item \textbf{Benar.} $3a+4b-12>0$.
\item \textbf{Salah.} $3a+4b=72\neq60$.
\item \textbf{Benar.} $3a+4b=72$.
\end{enumerate}
\end{pembox}

\begin{soal}{56}
$\displaystyle\int_5^9 f(x)\,dx=16$. Nilai
$\displaystyle\int_4^6 f(2x-3)\,dx$ adalah \ldots\hfill
Jawaban:\enspace\underline{\hspace{3cm}}
\end{soal}
\begin{pembox}
\jawaban{$8$}
Substitusi $u=2x-3$, $du=2\,dx$, $dx=\tfrac12du$.
Batas: $x=4\to u=5$; $x=6\to u=9$.
$\displaystyle\int_4^6 f(2x-3)\,dx=\tfrac12\int_5^9 f(u)\,du=\tfrac12\times16=8$.
\end{pembox}

\begin{soal}{57}
$AB$ diameter lingkaran, $A=(10,17)$, $B=(-2,1)$. Jari-jari lingkaran adalah \ldots
\begin{pilB}\item $20$\item $15$\item $10$\item $8$\item $4$\end{pilB}
\end{soal}
\begin{pembox}
\jawaban{(C) $r=10$}
$|AB|=\sqrt{(10-(-2))^2+(17-1)^2}=\sqrt{144+256}=\sqrt{400}=20$.
Jari-jari $=|AB|/2=10$.
\end{pembox}

\begin{soal}{58}
$f(x)=\log_3 x$.
Tentukan kebenaran setiap pernyataan berikut!

\par\smallskip
{\renewcommand{\arraystretch}{1.4}%
\begin{tabular}{@{}p{0.76\linewidth}cc@{}}
\hline\textbf{Pernyataan}&\textbf{Benar}&\textbf{Salah}\\\hline
(1)~$f(2x)+f\!\left(\dfrac{2}{x}\right)=\log_3\!\left(2x\cdot\dfrac{2}{x}\right)$.&$\bigcirc$&$\bigcirc$\\
(2)~$\log_3\!\left(2x\cdot\dfrac{2}{x}\right)=\log_3 4=f(4)$.&$\bigcirc$&$\bigcirc$\\
(3)~$f(2x)+f\!\left(\dfrac{2}{x}\right)=f(x^2)$.&$\bigcirc$&$\bigcirc$\\
(4)~$f(2x)+f\!\left(\dfrac{2}{x}\right)=f(4)$ (tidak bergantung $x$).&$\bigcirc$&$\bigcirc$\\\hline
\end{tabular}}
\end{soal}
\begin{pembox}
\jawaban{(1)~\textbf{B}\quad(2)~\textbf{B}\quad(3)~\textbf{S}\quad(4)~\textbf{B}}
$f(2x)+f\!\left(\tfrac2x\right)=\log_3(2x)+\log_3\!\left(\tfrac2x\right)=\log_3\!\left(2x\cdot\tfrac2x\right)=\log_3 4$.
\begin{enumerate}[leftmargin=*,topsep=2pt,itemsep=2pt]
\item \textbf{Benar.} Sifat $\log a+\log b=\log(ab)$.
\item \textbf{Benar.} $\log_3 4=f(4)$.
\item \textbf{Salah.} $f(x^2)=2\log_3 x$, bergantung pada $x$.
\item \textbf{Benar.} Hasilnya $f(4)$, konstan.
\end{enumerate}
\end{pembox}

\begin{soal}{59}
Data $D$: 5 bilangan bulat positif, tiga di antaranya $4,5,8$.
Rata-rata semua nilai median yang mungkin adalah \ldots
\begin{pilB}\item $5$\item $5\tfrac12$\item $6$\item $6\tfrac12$\item $7$\end{pilB}
\end{soal}
\begin{pembox}
\jawaban{(C) $6$}
Dua bilangan lain $p\leq q$ (positif). Setelah diurutkan 5 data, median adalah data ke-3.
Nilai median yang mungkin bergantung pada $p,q$:
median bisa $4,5,6,7,8$ (dengan memilih $p,q$ yang sesuai).
Rata-rata semua nilai median $=\dfrac{4+5+6+7+8}{5}=6$.
\end{pembox}

\begin{soal}{60}
Tentukan daerah asal fungsi $f(x)=\sqrt{|3x-7|-8}$.\\[0.4em]
Jawaban:\enspace\underline{\hspace{7cm}}
\end{soal}
\begin{pembox}
\jawaban{$x\leq-\dfrac{1}{3}$ atau $x\geq5$}
Syarat radikan $\geq0$: $|3x-7|-8\geq0\Rightarrow|3x-7|\geq8$.
Kasus 1: $3x-7\geq8\Rightarrow x\geq5$.
Kasus 2: $3x-7\leq-8\Rightarrow3x\leq-1\Rightarrow x\leq-\dfrac{1}{3}$.
Daerah asal: $\left\{x\;\middle|\;x\leq-\dfrac{1}{3}\text{ atau }x\geq5\right\}$.
\end{pembox}

\end{document}
